\documentclass[a4paper,12pt]{article}

\usepackage[french]{babel}
\usepackage[utf8]{inputenc}
\usepackage[T1]{fontenc}
% Sans fonte vectorielle en T1, pdfLaTeX fabrique des bitmaps (Type 3) : le PDF
% est flou au zoom et son texte mal copiable. Latin Modern corrige cela et se
% trouve dans toute distribution récente (MiKTeX, TeX Live) ; là où le paquet
% manque, cm-super joue le même rôle, et son absence ne doit pas empêcher la
% compilation — d'où le test plutôt qu'un \usepackage sec.
\IfFileExists{lmodern.sty}{\usepackage{lmodern}}{}
\usepackage{xcolor}
\usepackage{listings}
\usepackage{geometry}
\usepackage{enumitem}
\usepackage{amssymb}
\usepackage{graphicx}
\usepackage{float}
\usepackage{eurosym}
\usepackage{tikz}
\usetikzlibrary{positioning, fit, backgrounds}

% Ce document est plein de chemins, de commandes et de noms de variables en
% \texttt. Sans ces deux paquets, LaTeX refuse de les couper et les laisse
% déborder dans la marge — jusqu'à cinq centimètres sur l'URL de HackTheBox.
% « htt » autorise la césure dans la fonte à chasse fixe ; xurl permet de
% couper une URL n'importe où plutôt qu'aux seuls traits d'union.
\usepackage[htt]{hyphenat}
\usepackage{xurl}
\urlstyle{tt}
% Un \texttt contenant un chemin, une adresse ou un identifiant pointé n'offre
% aucun point de coupure : LaTeX le laisse alors déborder dans la marge. \code
% s'affiche comme \texttt mais se coupe où il faut. Son argument est pris tel
% quel : on y écrit _ et # sans les échapper.
\DeclareUrlCommand\code{\urlstyle{tt}}

\geometry{margin=2.5cm}

\lstdefinestyle{terminal}{
    backgroundcolor=\color{black},
    basicstyle=\ttfamily\small\color{white},
    frame=single,
    rulecolor=\color{gray},
    breaklines=true,
    columns=fullflexible
}

\lstdefinestyle{wincmd}{
    backgroundcolor=\color{black!85!blue},
    basicstyle=\ttfamily\small\color{white},
    frame=single,
    rulecolor=\color{gray},
    breaklines=true,
    columns=fullflexible
}

\lstset{style=terminal}

\newcommand{\screenshot}[2]{%
  \begin{figure}[H]
    \centering
    \fbox{\includegraphics[width=0.85\textwidth]{#1}}
    \caption{#2}
  \end{figure}
}

\newcommand{\screenshotplaceholder}[1]{%
  \begin{figure}[H]
    \centering
    \fbox{\begin{minipage}[c][5cm][c]{0.85\textwidth}
      \centering
      \color{gray}\itshape
      [Emplacement pour capture d'écran :\\ #1]
    \end{minipage}}
    \caption{#1}
  \end{figure}
}

% Encadré d'avertissement
\newcommand{\attention}[1]{%
  \begin{figure}[H]
  \noindent\fcolorbox{red!60!black}{red!6}{%
    \begin{minipage}{0.96\textwidth}
      \vspace{2pt}\textbf{\color{red!60!black}Attention —} #1\vspace{2pt}
    \end{minipage}}
  \end{figure}
}

\title{Lab « HackTheBox-like » sur GCP \\ avec une seule VM gratuite \\
       \large Passerelle VPN et machine cible fusionnées \\
       \large Pilotage depuis la plateforme cyberMans}
\author{Emmanuel Clement PREIRA}
\date{\today}

\begin{document}

\maketitle
\tableofcontents
\newpage

\section{Contexte}
\label{sec:contexte}

\begin{itemize}
    \item Plateforme : \textbf{Google Cloud Platform} (couche gratuite)
    \item Instance unique : \textbf{e2-micro} en \texttt{us-central1}
    \item Rôle 1 : \textbf{passerelle VPN} (serveur OpenVPN sur UDP 1194)
    \item Rôle 2 : \textbf{machine cible vulnérable} (DVWA dans Docker sur le port 80)
    \item Poste client : \textbf{Windows} avec \texttt{cmd.exe}
    \item \textbf{Pilotage} : la VM se démarre et s'éteint depuis cyberMans, exposée par ngrok
    \item Objectif : reproduire le mécanisme de HackTheBox --- un \texttt{.ovpn} qui fait entrer dans un réseau privé, et une cible joignable \textbf{uniquement} par le tunnel
    \item Résultat attendu : DVWA accessible à \texttt{http://10.8.0.1} une fois le tunnel monté
\end{itemize}

\subsection{Ce qui change par rapport à la première version}
\label{sec:changements}

Ce document corrige six points qui empêchaient le lab de fonctionner --- quatre
dans son branchement sur la plateforme, deux dans l'installation même du serveur
VPN --- et deux qui relevaient de la sécurité.

\begin{enumerate}
    \item \textbf{IP externe statique} (§\ref{sec:ip-pourquoi}). Sans réservation, chaque arrêt
          puis redémarrage change l'IP publique et \emph{invalide le
          \texttt{.ovpn} déjà distribué}. Or éteindre et rallumer est
          précisément ce que fait le bouton de la plateforme : sans IP fixe,
          le lab casse au premier cycle.
    \item \textbf{Adresse annoncée aux apprenants} (§\ref{sec:env}). Compute Engine
          rapporte l'adresse de la carte réseau dans le VPC
          (\code{10.128.0.2}). Le tunnel, lui, ne route que
          \code{10.8.0.0/24} : c'est \texttt{10.8.0.1} qu'il faut montrer.
          Sans ce réglage, la plateforme afficherait une adresse qui ne répond
          pas.
    \item \textbf{Réseau du lab déclaré à la plateforme} (§\ref{sec:env}). Il vaut
          \code{10.8.0.0/24}, pas la valeur par défaut \texttt{10.10.10.0/24}.
    \item \textbf{Compte de service GCP} (§\ref{sec:sa}), absent de la première
          version : sans lui, la plateforme ne peut rien démarrer.
    \item \textbf{Image DVWA} (§\ref{sec:docker}). \texttt{vulnerables/web-dvwa} n'est plus
          maintenue depuis des années ; l'image officielle actuelle est
          \texttt{ghcr.io/digininja/dvwa}.
    \item \textbf{Exposition du port 80} (§\ref{sec:barriere}). \texttt{-p 80:80} fait écouter
          DVWA sur \emph{toutes} les interfaces, IP publique comprise. Seul le
          pare-feu GCP l'en empêche : une seule barrière devant une application
          volontairement vulnérable. On en ajoute une seconde, locale.
    \item \textbf{Installation d'OpenVPN} (§\ref{sec:openvpn-install}). Le
          script d'angristan a été réécrit : il attend maintenant une
          sous-commande. \code{sudo ./openvpn-install.sh} seul affiche son aide
          et n'installe rien, sans erreur ni code de sortie non nul --- on croit
          avoir installé le serveur, et tout le reste du lab échoue.
    \item \textbf{Emplacement du \texttt{.ovpn}} (§\ref{sec:ovpn}). Lancé par
          \texttt{sudo}, le script écrit le profil dans le répertoire personnel
          de l'utilisateur, pas dans \texttt{/root/}. La commande de
          récupération de la première version échouait donc sur un fichier
          inexistant.
\end{enumerate}

\subsection{Convention de lecture}
\label{sec:convention}

\begin{itemize}
    \item Les blocs à fond noir sont des commandes \textbf{Linux}, dans la VM.
    \item Les blocs à fond bleu sombre sont des commandes \texttt{cmd.exe}, \textbf{sous Windows}.
    \item Les identifiants (ID de facturation, IP publique) sont à remplacer par les vôtres.
\end{itemize}

\subsection{Le modèle HackTheBox, et ce qu'on en reprend}
\label{sec:htb}

La documentation d'accès aux labs de HackTheBox (\emph{Introduction to Lab
Access}, \url{https://help.hackthebox.com/en/articles/5185687-introduction-to-lab-access}) commence par écarter une
confusion qui vaut d'être reprise telle quelle devant une promotion : son VPN
n'est pas un VPN de confidentialité. Un VPN personnel chiffre le trafic et le
fait ressortir ailleurs, pour masquer une activité ou une localisation. Celui
de HackTheBox est de l'autre famille, celle des VPN d'entreprise : il sert à
\textbf{entrer} dans un réseau interne depuis Internet, comme un salarié à
distance atteint les ressources de son employeur.

De là découle la phrase qui résume tout le mécanisme : le profil
« \emph{vous placera dans le même sous-réseau IP que les machines
vulnérables, vous permettant de les contacter (et de les attaquer)} ». Rien
d'autre n'ouvre la porte : pas de redirection de port, pas d'adresse publique
sur la cible. Le tunnel \emph{est} l'accès.

Ce que l'article décrit ensuite, ce sont les réglages du téléchargement :

\begin{itemize}
    \item un bouton \textbf{Connecter à HTB}, présent sur toutes les vues du
          tableau de bord, à côté de l'avatar ;
    \item le \textbf{type de contenu} visé (Machines, Starting Point,
          Fortresses, Pro Labs, Seasonal, Release Arena) ;
    \item le \textbf{niveau} d'accès (gratuit ou VIP+) et la \textbf{région}
          (Europe, États-Unis, Australie, Singapour) ;
    \item le \textbf{serveur} précis, choisi à la main --- et non attribué au
          hasard, pour que des coéquipiers puissent se retrouver sur le même ;
    \item le \textbf{protocole}, TCP ou UDP ;
    \item le bouton de \textbf{téléchargement} du pack \texttt{.ovpn} ;
    \item la \textbf{régénération}, obtenue en basculant d'un serveur à
          l'autre quand le fichier est mal formé.
\end{itemize}

\attention{L'article ne parle \textbf{pas} du démarrage des machines : il ne
traite que de l'accès au réseau. Le bouton qui allume une cible vit sur la
page de la machine, ailleurs dans l'interface. Ce sont deux mécanismes
distincts, et c'est bien ainsi qu'ils sont séparés dans cyberMans --- page
\textbf{Accès VPN} d'un côté, page \textbf{Machines} de l'autre.}

\subsubsection*{Où chaque geste se retrouve dans cyberMans}

\begin{description}[leftmargin=0cm, style=nextline]
  \item[\textbf{Le téléchargement du profil}] C'est la page \textbf{Accès
        VPN}. En mode \texttt{generated}, la plateforme reprend presque à
        l'identique les menus de HackTheBox : choix du protocole --- avec les
        mêmes indications, « plus rapide » pour UDP, « si votre réseau bloque
        l'UDP » pour TCP --- et régénération du profil en un clic. Le présent
        lab utilise l'autre mode, \texttt{uploaded} (§\ref{sec:depot}) :
        l'administrateur dépose un seul fichier, le même pour tous, et ces deux
        réglages disparaissent de l'écran puisqu'il n'y a plus rien à choisir
        ni à réémettre.
  \item[\textbf{Le choix du serveur et de la région}] Sans objet : il n'y a
        qu'une passerelle, celle de la VM.
  \item[\textbf{Le démarrage de la cible}] Page \textbf{Machines} :
        \texttt{POST /api/machine/start}, \texttt{POST /api/machine/stop}, et
        \texttt{GET /api/machine/status} qui rend l'état et, quand la machine
        tourne, son adresse (§\ref{sec:env}).
\end{description}

\subsubsection*{Trois écarts, et ce qu'ils coûtent}

\begin{description}[leftmargin=0cm, style=nextline]
  \item[\textbf{La cible est partagée}] C'est l'écart le plus lourd, et il
        vaut la peine de noter que HackTheBox l'a lui-même jugé intenable :
        depuis le 15 janvier 2026, les comptes gratuits et VIP reçoivent une
        \emph{instance dédiée} par machine, précisément pour que « personne
        n'attaque la même machine que vous ». Ici, tous les apprenants
        attaquent la même DVWA : ils se marchent sur les pieds, et celui qui
        clique sur « Arrêter » éteint la machine de tout le monde. Acceptable
        pour une démonstration ou un lab d'équipe, pas pour une promotion.

        La plateforme sait désormais faire autrement : avec
        \texttt{APP\_GCP\_TARGET\_MODE=per-user}, elle \emph{crée} une
        instance par apprenant au démarrage et la \emph{détruit} à l'arrêt,
        comme HackTheBox. Ce lab-ci ne l'emploie pas, et pour une seule
        raison : chaque apprenant actif devient une machine facturée, ce qui
        sort de la couche gratuite dès le second. Le compromis est détaillé
        dans \texttt{docs/GCP.md}.
  \item[\textbf{Le profil est collectif}] Chez HackTheBox le pack est
        personnel. Ici un seul fichier circule : on ne peut pas révoquer un
        apprenant sans redéposer un profil et le redistribuer à tous
        (\texttt{docs/VPN.md}). Le mode \texttt{generated} rend cette
        révocation individuelle, au prix d'une infrastructure de certificats.
  \item[\textbf{La passerelle et la cible sont la même machine}]
        (§\ref{sec:ip-pourquoi}) Chez HackTheBox elles sont distinctes ; ici
        l'adresse de la cible est celle du serveur VPN lui-même,
        \texttt{10.8.0.1}. C'est ce qui rend
        \texttt{APP\_MACHINE\_ADDRESS} nécessaire (§\ref{sec:env}).
\end{description}

\attention{Le profil déposé ici est en \textbf{UDP 1194} seulement. HackTheBox
propose TCP pour les réseaux qui perdent des paquets ou qui filtrent l'UDP ---
cas courant sur un Wi-Fi d'entreprise ou d'école. Un apprenant dans ce cas ne
montera pas le tunnel, et aucun message de la plateforme ne le lui dira. Si la
promotion est sur un réseau filtré, relancez le script d'installation
(§\ref{sec:openvpn-install}) en choisissant TCP, et déposez ce second profil à
la place.}

\section{Architecture}
\label{sec:archi}

\subsection{Le principe : une seule VM, deux rôles}
\label{sec:principe}

L'unique VM joue \textbf{deux rôles en même temps} :

\begin{enumerate}
    \item \textbf{Porte d'entrée VPN} --- serveur OpenVPN. Son interface de
          tunnel (\texttt{tun0}) porte l'adresse \texttt{10.8.0.1}.
    \item \textbf{Cible vulnérable} --- DVWA dans Docker, port 80 de cette
          \emph{même} machine.
\end{enumerate}

On attaque donc directement \texttt{10.8.0.1} : l'adresse de la VM \emph{vue de
l'intérieur du tunnel}.

\begin{figure}[H]
\centering
\begin{tikzpicture}[
    font=\small,
    box/.style={rounded corners=3pt, draw, thick, align=center,
                minimum width=3.4cm, minimum height=1.1cm},
    pc/.style={box, fill=gray!15},
    vpn/.style={box, fill=teal!15},
    target/.style={box, fill=orange!18},
    vm/.style={rounded corners=8pt, draw, thick, dashed, inner sep=14pt}
]
\node[pc] (pc) {\textbf{PC Windows}\\\texttt{tun0 = 10.8.0.2}};
\node[vpn, right=5cm of pc, yshift=10mm] (ovpn)
    {\textbf{Serveur OpenVPN}\\\texttt{tun0 = 10.8.0.1}};
\node[target, below=10mm of ovpn] (dvwa)
    {\textbf{DVWA} (Docker :80)\\cible vulnérable};
\begin{scope}[on background layer]
\node[vm, fit=(ovpn)(dvwa),
      label={[align=center]above:\textbf{VM e2-micro}\\\footnotesize IP publique \emph{statique}}] (vm) {};
\end{scope}
\draw[->, thick] (pc) -- node[above, align=center]
    {\footnotesize tunnel chiffré\\\footnotesize UDP 1194} (vm.west);
\end{tikzpicture}
\caption{Passerelle VPN et cible fusionnées. On attaque \texttt{10.8.0.1} à
         travers le tunnel ; le port 80 n'est jamais exposé à Internet.}
\end{figure}

\subsection{Le trajet d'un paquet}
\label{sec:trajet}

Quand on tape \texttt{http://10.8.0.1} dans le navigateur Windows :

\begin{enumerate}
    \item Le PC fabrique \texttt{[10.8.0.2 $\rightarrow$ 10.8.0.1:80]} ; sa table de routage envoie \code{10.8.0.0/24} vers l'interface VPN.
    \item OpenVPN \textbf{chiffre} et emballe dans un UDP ordinaire \texttt{[IP\_PC $\rightarrow$ IP\_publique:1194]}.
    \item La VM \textbf{déchiffre} et retrouve le paquet sur \texttt{tun0}.
    \item \texttt{10.8.0.1} est \textbf{sa propre adresse} : livraison directe au conteneur. \textbf{Aucun routage, aucun NAT.}
    \item DVWA répond par le même tunnel.
\end{enumerate}

Votre analyse de la première version était juste : \textbf{pas de
\texttt{MASQUERADE}, pas de \texttt{--can-ip-forward}}, et
\texttt{net.ipv4.ip\_forward=1} est \textbf{inoffensif mais inutile} ici. Le
forwarding ne sert qu'à router vers d'\emph{autres} hôtes ; la cible \emph{est}
la passerelle.

\subsection{Pourquoi la cible reste cachée}
\label{sec:cachee}

Le pare-feu GCP n'ouvre que \textbf{UDP 1194} et \textbf{TCP 22}. Le port 80
n'est jamais ouvert publiquement : un scan de l'IP publique ne révèle que le
VPN et SSH.

\attention{Cette propriété ne tient qu'à la règle de pare-feu. DVWA écoute en
réalité sur \texttt{0.0.0.0:80}, donc une règle ajoutée par mégarde, un projet
mal configuré ou une VM déplacée dans un autre réseau exposerait à Internet une
application \emph{conçue} pour être vulnérable. La section~\ref{sec:barriere} ajoute une seconde
barrière, locale à la VM, qui ne dépend d'aucun réglage GCP.}

\subsection{Le compromis}
\label{sec:compromis}

Une seule machine, donc \textbf{pas de pivoting} ni de mouvement latéral. Pour
du pentest web (DVWA, Juice Shop) c'est parfait ; pour le déplacement en réseau
interne il faudra deux VM.

\section{Installation de gcloud CLI sur Windows}
\label{sec:gcloud}

\begin{lstlisting}[style=wincmd]
powershell -Command "(New-Object Net.WebClient).DownloadFile('https://dl.google.com/dl/cloudsdk/channels/rapid/GoogleCloudSDKInstaller.exe', '%TEMP%\GoogleCloudSDKInstaller.exe')"
%TEMP%\GoogleCloudSDKInstaller.exe
\end{lstlisting}

Options : \textbf{Single User}, cocher \textbf{Core Libraries}, décocher
« Bundled Python » si Python est déjà installé.

Vérification dans un \textbf{nouveau} \texttt{cmd.exe} :

\begin{lstlisting}[style=wincmd]
gcloud --version
\end{lstlisting}

\screenshotplaceholder{cmd.exe : sortie de \texttt{gcloud --version}}

\section{Création du projet GCP}
\label{sec:projet}

\begin{lstlisting}[style=wincmd]
gcloud projects create lab-htb-free-2026 --name="Lab HTB Free"
gcloud config set project lab-htb-free-2026
gcloud config set compute/zone us-central1-a
gcloud config set compute/region us-central1
gcloud config list
\end{lstlisting}

Liaison au compte de facturation (le préfixe \texttt{beta} n'est plus nécessaire) :

\begin{lstlisting}[style=wincmd]
gcloud billing accounts list
gcloud billing projects link lab-htb-free-2026 --billing-account=012345-6789AB-CDEF01
\end{lstlisting}

Activation de l'API Compute (30 s à 2 min) :

\begin{lstlisting}[style=wincmd]
gcloud services enable compute.googleapis.com
\end{lstlisting}

\screenshotplaceholder{cmd.exe : projet actif et facturation liée}

\section{Adresse IP statique}
\label{sec:ip}

\subsection{Pourquoi c'est obligatoire ici}
\label{sec:ip-pourquoi}

Une IP \emph{éphémère} est libérée à chaque arrêt de la VM et change au
redémarrage. Le \texttt{.ovpn} contient la ligne \texttt{remote <IP> 1194} :
une IP qui change \textbf{invalide tous les profils déjà distribués}.

Comme le bouton de la plateforme éteint et rallume la VM, le lab casserait dès
le premier cycle. Il faut donc une \textbf{IP statique réservée}.

\subsection{Réservation}
\label{sec:ip-reservation}

\begin{lstlisting}[style=wincmd]
gcloud compute addresses create lab-htb-ip --region=us-central1
gcloud compute addresses describe lab-htb-ip --region=us-central1 --format="get(address)"
\end{lstlisting}

Notez l'adresse : c'est elle qui ira dans le \texttt{.ovpn} et nulle part ailleurs.

\attention{Une IP statique est facturée qu'elle soit utilisée ou non (environ
3~\$/mois, et davantage quand elle est réservée sans être attachée). C'est le
prix à payer pour un \texttt{.ovpn} qui survit aux redémarrages. Si vous
renoncez à piloter la VM depuis la plateforme, une IP éphémère suffit --- mais
il faudra rediffuser le profil après chaque démarrage.}

\section{Création de la VM}
\label{sec:vm}

Trois critères pour rester dans la couche gratuite : zone
\texttt{us-central1} / \texttt{us-west1} / \texttt{us-east1}, type
\texttt{e2-micro}, disque 30~Go standard.

\textbf{Une seule ligne} (le \texttt{\^} de continuation est capricieux sous \texttt{cmd.exe}) :

\begin{lstlisting}[style=wincmd]
gcloud compute instances create lab-htb-single --zone=us-central1-a --machine-type=e2-micro --image-family=debian-12 --image-project=debian-cloud --boot-disk-size=30GB --boot-disk-type=pd-standard --tags=vpn-server --address=lab-htb-ip
\end{lstlisting}

Le \texttt{--address=lab-htb-ip} attache l'IP réservée : c'est la seule
différence avec la première version, et elle est décisive.

\screenshotplaceholder{cmd.exe : création de la VM avec son IP statique}

\section{Pare-feu}
\label{sec:pare-feu}

\begin{lstlisting}[style=wincmd]
gcloud compute firewall-rules create allow-openvpn --allow=udp:1194 --source-ranges=0.0.0.0/0 --target-tags=vpn-server
\end{lstlisting}

\begin{lstlisting}[style=wincmd]
gcloud compute firewall-rules create allow-ssh --allow=tcp:22 --source-ranges=0.0.0.0/0
\end{lstlisting}

Le port 80 n'est volontairement \emph{pas} ouvert.

\screenshotplaceholder{Console GCP : les deux règles de pare-feu}

\section{OpenVPN}
\label{sec:openvpn}

\subsection{Connexion}
\label{sec:ssh}

\begin{lstlisting}[style=wincmd]
gcloud compute ssh lab-htb-single --zone=us-central1-a
\end{lstlisting}

\subsection{Installation}
\label{sec:openvpn-install}

\begin{lstlisting}
curl -O https://raw.githubusercontent.com/angristan/openvpn-install/master/openvpn-install.sh
chmod +x openvpn-install.sh
\end{lstlisting}

\attention{Le script attend désormais une \textbf{sous-commande}. Lancé seul,
\code{sudo ./openvpn-install.sh} se contente d'afficher son aide et
\textbf{n'installe rien} --- en silence, et avec un code de sortie 0, donc sans
rien qui ressemble à une erreur. C'est le piège : on croit avoir installé
OpenVPN, et la suite du lab échoue sans raison apparente.}

Deux manières de faire. La première ne pose aucune question et se rejoue à
l'identique :

\begin{lstlisting}
sudo ./openvpn-install.sh install \
    --endpoint VOTRE.IP.STATIQUE \
    --port 1194 --protocol udp \
    --no-client-ipv6 \
    --no-route-internet \
    --client lab-htb
\end{lstlisting}

\begin{description}[leftmargin=0cm, style=nextline]
  \item[\code{--endpoint}] L'adresse écrite dans le \texttt{.ovpn}. C'est elle
        qui règle le problème du NAT 1:1 : la donner explicitement évite que le
        script détecte l'adresse interne (\code{10.128.0.2}) et produise un
        profil injoignable depuis Internet. Mettez l'\textbf{IP statique} de la
        section~\ref{sec:ip-reservation}.
  \item[\code{--no-route-internet}] Seul le réseau du lab passe par le tunnel.
        Sans cette option, \emph{tout} le trafic de l'apprenant --- sa
        navigation comprise --- sort par la VM : vous payez l'egress GCP, et une
        e2-micro devient le goulot d'étranglement de la promotion. C'est aussi
        ce que fait HackTheBox, qui ne route que ses sous-réseaux de lab.
  \item[\code{--client lab-htb}] Le nom du premier profil, donc le nom du
        fichier : \texttt{lab-htb.ovpn}.
\end{description}

La seconde relance le questionnaire d'autrefois, utile pour voir les choix
défiler :

\begin{lstlisting}
sudo ./openvpn-install.sh interactive
\end{lstlisting}

\subsection{Réponses au questionnaire}
\label{sec:openvpn-reponses}

Uniquement si vous avez choisi \code{interactive} --- la forme non interactive
ci-dessus ne demande rien.

\begin{itemize}
    \item \textbf{IP address} : l'\textbf{IP statique} de la section~\ref{sec:ip-reservation} --- impératif
    \item \textbf{IPv6} : \texttt{n} \quad\textbullet\quad \textbf{Port} : \texttt{1194} \quad\textbullet\quad \textbf{Protocol} : \texttt{UDP}
    \item \textbf{DNS} : au choix (ex. \texttt{1.1.1.1})
    \item \textbf{Client name} : \texttt{lab-htb}
\end{itemize}

\attention{À cause du NAT 1:1 de GCP, le script propose l'IP \textbf{interne}
(\code{10.128.0.2}). Saisissez l'IP publique \textbf{à la main}, sinon le
\texttt{.ovpn} pointera vers une adresse injoignable depuis Internet. C'est
exactement ce que \code{--endpoint} évite dans la forme non interactive.}

\screenshotplaceholder{Terminal VM : questionnaire du script OpenVPN}

\section{DVWA sur la même VM}
\label{sec:dvwa}

\subsection{Swap}
\label{sec:swap}

\texttt{e2-micro} n'a qu'1~Go de RAM.

\begin{lstlisting}
sudo fallocate -l 2G /swapfile || sudo dd if=/dev/zero of=/swapfile bs=1M count=2048
sudo chmod 600 /swapfile
sudo mkswap /swapfile
sudo swapon /swapfile
echo '/swapfile none swap sw 0 0' | sudo tee -a /etc/fstab
free -h
\end{lstlisting}

Le repli en \texttt{dd} couvre les systèmes de fichiers où \texttt{fallocate} échoue.

\subsection{Docker et DVWA}
\label{sec:docker}

\begin{lstlisting}
sudo apt update && sudo apt install -y docker.io
sudo systemctl enable --now docker
sudo docker run -d --name dvwa --restart=always -p 80:80 ghcr.io/digininja/dvwa:latest
sudo docker ps
\end{lstlisting}

\attention{L'image \texttt{vulnerables/web-dvwa} de la première version n'est
plus maintenue depuis des années et traîne des vulnérabilités de base qui ne
font pas partie de l'exercice. \texttt{ghcr.io/digininja/dvwa} est l'image
officielle du projet, toujours suivie.}

\subsection{Seconde barrière devant le port 80}
\label{sec:barriere}

\texttt{-p 80:80} publie sur \emph{toutes} les interfaces. On restreint
localement le port 80 au seul tunnel :

\begin{lstlisting}
sudo apt install -y iptables-persistent
sudo iptables -I DOCKER-USER -i tun0 -p tcp --dport 80 -j ACCEPT
sudo iptables -A DOCKER-USER -p tcp --dport 80 -j DROP
sudo netfilter-persistent save
\end{lstlisting}

La chaîne \texttt{DOCKER-USER} est consultée \emph{avant} les règles que Docker
génère : c'est le seul endroit où une règle survit à un redémarrage du
conteneur. Désormais, même si la règle de pare-feu GCP venait à sauter, le port
80 ne répondrait qu'aux paquets arrivant par \texttt{tun0}.

\screenshotplaceholder{Terminal VM : \texttt{docker ps} et \texttt{free -h}}

\section{Récupération du \texttt{.ovpn}}
\label{sec:ovpn}

Le script annonce le chemin en dernière ligne :
\texttt{The configuration file has been written to \ldots}. Lisez-la, elle est
la source de vérité.

Lancé avec \texttt{sudo} depuis une session \texttt{gcloud compute ssh}, le
script écrit le profil dans \textbf{votre} répertoire personnel --- il suit
\texttt{SUDO\_USER} --- et non dans \texttt{/root/}. Il vous en donne aussi la
propriété. Un seul transfert suffit donc, sans \texttt{sudo} ni passage par
\texttt{/tmp} :

\begin{lstlisting}[style=wincmd]
gcloud compute scp lab-htb-single:~/lab-htb.ovpn ./lab-htb.ovpn --zone=us-central1-a
\end{lstlisting}

\attention{Si vous avez lancé le script depuis un \emph{shell} déjà root
(\texttt{sudo -i} puis \texttt{./openvpn-install.sh \ldots}), alors
\texttt{SUDO\_USER} vaut \texttt{root} et le profil est bien dans
\texttt{/root/}. Il faut alors le sortir avant de le transférer :
\texttt{sudo cp /root/lab-htb.ovpn /tmp/ \&\& sudo chmod 644
/tmp/lab-htb.ovpn}.}

\subsection{Vérification du contenu}
\label{sec:ovpn-verif}

Ouvrez le fichier dans un éditeur et contrôlez la ligne \texttt{remote} :

\begin{lstlisting}[style=wincmd]
findstr remote lab-htb.ovpn
\end{lstlisting}

Elle doit porter l'\textbf{IP statique}, jamais \texttt{10.128.0.x}.

\section{Connexion depuis Windows}
\label{sec:windows}

\begin{itemize}
    \item OpenVPN GUI $\rightarrow$ \textbf{Import file\ldots} $\rightarrow$ \texttt{lab-htb.ovpn}
    \item \textbf{Lancer OpenVPN GUI en administrateur} : sans cela la route vers \code{10.8.0.0/24} n'est pas ajoutée et le \texttt{ping} échoue malgré un tunnel « connecté ».
    \item Clic droit $\rightarrow$ \textbf{Connect}. Attendre \texttt{Initialization Sequence Completed}.
\end{itemize}

Vérification dans un second \texttt{cmd.exe} :

\begin{lstlisting}[style=wincmd]
ipconfig | findstr /i "tun"
route print | findstr "10.8.0"
ping 10.8.0.1
\end{lstlisting}

\screenshotplaceholder{cmd.exe : \texttt{ping 10.8.0.1} réussi}

\attention{Le tunnel doit \textbf{rester ouvert} pendant toute la session.
Fermer OpenVPN GUI --- ou, sous Linux, la fenêtre où tourne \texttt{openvpn} ---
tue le processus et coupe l'accès au lab sur-le-champ. L'icône dans la zone de
notification doit rester verte.}

\attention{\textbf{Un seul VPN à la fois.} Si l'apprenant a déjà un VPN
personnel actif (celui de son école, un service de confidentialité), les deux
jeux de routes entrent en conflit et le tunnel du lab monte sans rien joindre.
C'est le même piège que signale HackTheBox entre son VPN et son Pwnbox. Il faut
couper l'autre avant de se connecter.}

\subsection{Depuis Linux (Kali, Parrot)}
\label{sec:linux-client}

Le profil est le même ; seul le lancement change. Dans un terminal :

\begin{lstlisting}[style=terminal]
sudo openvpn lab-htb.ovpn
\end{lstlisting}

Attendre la ligne \texttt{Initialization Sequence Completed}, puis
\textbf{laisser ce terminal ouvert} et en ouvrir un second pour travailler.
C'est la marche à suivre que décrit la documentation de HackTheBox, et
\texttt{sudo} y est indispensable pour la même raison qu'« administrateur »
sous Windows : sans lui, la route vers \code{10.8.0.0/24} n'est pas posée.

\section{Attaque de DVWA}
\label{sec:attaque}

\begin{lstlisting}[style=wincmd]
start http://10.8.0.1
\end{lstlisting}

Identifiants : \texttt{admin} / \texttt{password}. Cliquer sur \textbf{Create /
Reset Database} au premier lancement.

\screenshotplaceholder{Navigateur : DVWA sur \texttt{http://10.8.0.1}}

\section{Branchement sur la plateforme cyberMans}
\label{sec:plateforme}

C'est la partie absente de la première version. L'objectif : que le bouton
« Démarrer / Arrêter » de la plateforme pilote cette VM, et que les apprenants
y téléchargent le \texttt{.ovpn}.

\subsection{Compte de service GCP}
\label{sec:sa}

La plateforme agit au nom d'un compte de service. Dans \texttt{cmd.exe} :

\begin{lstlisting}[style=wincmd]
gcloud iam service-accounts create cybermans-lab --display-name="cyberMans - pilotage du lab"
\end{lstlisting}

\begin{lstlisting}[style=wincmd]
gcloud projects add-iam-policy-binding lab-htb-free-2026 --member="serviceAccount:cybermans-lab@lab-htb-free-2026.iam.gserviceaccount.com" --role="roles/compute.instanceAdmin.v1"
\end{lstlisting}

\begin{lstlisting}[style=wincmd]
gcloud iam service-accounts keys create cle-lab.json --iam-account=cybermans-lab@lab-htb-free-2026.iam.gserviceaccount.com
\end{lstlisting}

\code{roles/compute.instanceAdmin.v1} suffit : il couvre
\code{instances.start}, \code{instances.stop} et \code{instances.get}. Ne
donnez pas \code{roles/editor}, qui ouvrirait le projet entier à un service
dont le seul travail est d'actionner un interrupteur.

\attention{\texttt{cle-lab.json} est un secret de même nature qu'un mot de
passe d'administration. Il reste sur le serveur qui héberge la plateforme,
n'entre jamais dans un dépôt Git, et n'est jamais servi au navigateur.}

\subsection{Dépôt du \texttt{.ovpn} dans la plateforme}
\label{sec:depot}

Connectez-vous à cyberMans en administrateur, puis
\textbf{Administration $\rightarrow$ Profil VPN}, et déposez
\texttt{lab-htb.ovpn}. Les apprenants le téléchargeront depuis
« Accès VPN ».

\begin{itemize}
    \item La plateforme vérifie que le fichier est bien un profil \emph{client}
          (directives \texttt{client} et \texttt{remote}).
    \item Elle peut afficher « \emph{Ce profil ne route pas le réseau du lab} ».
          \textbf{C'est normal ici} : OpenVPN installe tout seul la route vers
          le sous-réseau du tunnel, sans qu'une ligne \texttt{route} figure dans
          le profil. L'avertissement ne bloque rien --- vérifiez simplement que
          \texttt{ping 10.8.0.1} répond, c'est la seule preuve qui compte.
\end{itemize}

\subsection{Variables d'environnement de la plateforme}
\label{sec:env}

\begin{lstlisting}[style=wincmd]
APP_MACHINE_PROVIDER=gcp
GCP_PROJECT_ID=lab-htb-free-2026
GCP_ZONE=us-central1-a
GCP_INSTANCE_NAME=lab-htb-single
GOOGLE_APPLICATION_CREDENTIALS=/chemin/vers/cle-lab.json

APP_MACHINE_ADDRESS=10.8.0.1
APP_VPN_ENABLED=true
APP_VPN_SOURCE=uploaded
APP_VPN_LAB_NETWORK=10.8.0.0/24
\end{lstlisting}

Les quatre dernières lignes sont celles qu'on oublie, et chacune produit une
panne silencieuse :

\begin{description}[leftmargin=2.2cm, style=nextline]
  \item[\texttt{APP\_VPN\_ENABLED}] \texttt{true} : sans cela la page
        \textbf{Accès VPN} reste fermée et personne ne télécharge le profil.
  \item[\texttt{APP\_MACHINE\_ADDRESS}] Compute Engine rapporte l'adresse de la
        carte réseau dans le VPC (\code{10.128.0.2}). Or le tunnel ne route
        que \code{10.8.0.0/24} : sans ce réglage, la plateforme afficherait
        \code{10.128.0.2}, qui ne répondrait pas. L'apprenant conclurait à une
        machine en panne.
  \item[\texttt{APP\_VPN\_SOURCE}] \texttt{uploaded} : le profil vient de votre
        serveur, pas de la plateforme. Avec \texttt{generated} (le défaut), elle
        tenterait d'émettre ses propres certificats.
  \item[\texttt{APP\_VPN\_LAB\_NETWORK}] Le réseau annoncé aux apprenants sur
        la page \textbf{Accès VPN}. C'est aussi lui que la plateforme cherche
        dans le \texttt{.ovpn} déposé : si le profil ne le route pas, l'écran
        d'administration le signale. La valeur par défaut est
        \texttt{10.10.10.0/24} ; votre tunnel est en \code{10.8.0.0/24}.
\end{description}

\subsection{Exposition par ngrok}
\label{sec:ngrok}

La plateforme se lance par \texttt{docker compose up -d} ; Nginx sert
l'application sur le port \texttt{3000} et relaie lui-même \texttt{/api} vers
le backend. Navigateur et API partagent donc la \emph{même} origine : il n'y a
rien à déclarer côté CORS, et c'est cette même origine qui permet au cookie de
session (\texttt{SameSite=Strict}) d'être envoyé.

\begin{lstlisting}[style=wincmd]
ngrok http 3000
\end{lstlisting}

ngrok rend une URL du type \texttt{https://xxxx.ngrok-free.app}. Deux variables
de la plateforme doivent en tenir compte :

\begin{lstlisting}[style=wincmd]
APP_PUBLIC_URL=https://xxxx.ngrok-free.app
APP_COOKIE_SECURE=true
\end{lstlisting}

\begin{description}[leftmargin=2.6cm, style=nextline]
  \item[\texttt{APP\_PUBLIC\_URL}] L'adresse écrite dans les liens que la
        plateforme émet --- celui de réinitialisation de mot de passe en
        particulier. Restée sur \texttt{http://localhost:3000}, elle produit des
        liens qui ne mènent nulle part depuis un autre poste.
  \item[\texttt{APP\_COOKIE\_SECURE}] \texttt{true} marque le cookie de session
        \texttt{Secure}, ce que l'on doit faire dès que le trafic passe en
        HTTPS : ngrok termine le TLS. Les deux valeurs « fonctionnent » dans un
        navigateur ; seule \texttt{true} empêche le cookie de repartir un jour
        en clair.
\end{description}

\attention{Avec un compte ngrok gratuit, cette URL \textbf{change à chaque
redémarrage} du tunnel. Il faut alors remettre \texttt{APP\_PUBLIC\_URL} à jour
et relancer la plateforme. Un domaine statique (offert sur le palier gratuit,
un seul par compte) évite ce va-et-vient.}

La première visite affiche la page d'avertissement de ngrok (« You are about to
visit\ldots ») : c'est le palier gratuit, pas une erreur. On clique sur
\textbf{Visit Site} et elle ne revient plus pour la session.

\attention{Le tunnel ngrok expose la plateforme, \textbf{pas} le lab. La cible
reste joignable uniquement par le \texttt{.ovpn} : l'URL ngrok sert à cliquer
sur « Démarrer », jamais à atteindre DVWA.}

\subsection{Ce que l'apprenant fait, dans l'ordre}
\label{sec:apprenant}

\begin{enumerate}
    \item Il ouvre l'URL ngrok et se connecte.
    \item \textbf{Accès VPN} $\rightarrow$ il télécharge \texttt{lab-htb.ovpn}, l'importe dans OpenVPN GUI (\textbf{en administrateur}) et se connecte.
    \item \textbf{Machines} $\rightarrow$ il clique sur \textbf{Démarrer la machine}. Le badge passe par \texttt{Démarrage en cours}, le bouton reste inactif, l'état se rafraîchit tout seul.
    \item Une fois \texttt{En marche}, l'adresse \texttt{10.8.0.1} s'affiche.
    \item Il ouvre \texttt{http://10.8.0.1} et travaille sur DVWA.
    \item En partant, il clique sur \textbf{Arrêter la machine}.
\end{enumerate}

\attention{La cible est \textbf{partagée} : tous les apprenants pilotent la même
VM, et le second à cliquer sur « Arrêter » éteint la machine du premier. Pour un
lab d'équipe ou une démonstration, c'est sans conséquence. Pour une promotion
entière, il faudra une VM par apprenant --- ce que ce montage à une seule
machine ne permet pas.}

\section{Ordre de démarrage après un arrêt}
\label{sec:ordre}

La plateforme rallume la VM, mais elle n'attend pas que les services soient
prêts. Comptez \textbf{30 à 60 secondes} entre l'état \texttt{En marche} et le
moment où DVWA répond : Docker doit relancer le conteneur, et OpenVPN
redémarrer.

Les deux services reviennent seuls : \texttt{--restart=always} pour le
conteneur, \texttt{systemd} pour OpenVPN. Vérification après un cycle :

\begin{lstlisting}
sudo systemctl status openvpn-server@server
sudo docker ps
sudo iptables -L DOCKER-USER -n
\end{lstlisting}

La troisième commande confirme que la règle de la section~\ref{sec:barriere} a bien survécu --- c'est
le rôle de \texttt{iptables-persistent}.

\section{Facturation}
\label{sec:facturation}

\begin{itemize}
    \item \textbf{VM \texttt{e2-micro}} : gratuite dans les zones éligibles, une par compte et par mois.
    \item \textbf{IP externe IPv4} : facturée depuis février 2024, qu'elle soit
          statique ou éphémère, dès lors qu'elle existe. Comptez environ
          3~\$/mois. Une IP statique \emph{réservée mais non attachée} coûte
          davantage : si vous supprimez la VM, libérez aussi l'adresse.
    \item \textbf{Disque} : 30~Go standard, dans la couche gratuite.
\end{itemize}

Libération de l'IP lors d'une destruction complète :

\begin{lstlisting}[style=wincmd]
gcloud compute instances delete lab-htb-single --zone=us-central1-a --quiet
gcloud compute addresses delete lab-htb-ip --region=us-central1 --quiet
gcloud compute firewall-rules delete allow-openvpn allow-ssh --quiet
\end{lstlisting}

Alerte de budget : \texttt{console.cloud.google.com/billing} $\rightarrow$
Budgets et alertes $\rightarrow$ 5~\$, seuils 50 / 90 / 100~\%.

\section{Dépannage}
\label{sec:depannage}

\begin{description}[leftmargin=0.6cm, style=nextline]
  \item[\textmd{\texttt{The required property [project] is not currently set}}]
        \texttt{gcloud config set project lab-htb-free-2026}.
  \item[\textmd{\texttt{--zone n'est pas reconnu en tant que commande}}]
        La commande multi-lignes a été coupée par \texttt{cmd.exe}. Tout sur une ligne.
  \item[\textmd{Pas de \texttt{Initialization Sequence Completed}}]
        Le \texttt{.ovpn} pointe vers l'IP interne. Corriger la ligne \texttt{remote}.
  \item[\textmd{Connecté, mais \texttt{ping 10.8.0.1} échoue}]
        OpenVPN GUI n'a pas été lancé en administrateur. Sinon, dans la VM :
        \texttt{sudo systemctl status openvpn-server@server}.
  \item[\textmd{\texttt{http://10.8.0.1} ne répond pas}]
        \texttt{sudo docker ps}, puis \code{sudo docker logs dvwa}. Vérifier
        aussi \code{sudo iptables -L DOCKER-USER -n} : une règle trop large
        bloquerait \texttt{tun0}.
  \item[\textmd{Le \texttt{.ovpn} ne marche plus après un redémarrage}]
        L'IP n'est pas statique (§\ref{sec:ip}). C'est la panne la plus fréquente dès que
        la plateforme pilote la VM.
  \item[\textmd{La plateforme affiche \code{10.128.0.2}}]
        \texttt{APP\_MACHINE\_ADDRESS} n'est pas renseigné (§\ref{sec:env}).
  \item[\textmd{Déconnecté à chaque ouverture de l'URL ngrok}]
        L'URL a changé depuis la dernière session : le cookie était posé pour
        l'ancien nom d'hôte et n'est plus envoyé. Il faut se reconnecter, et
        remettre \texttt{APP\_PUBLIC\_URL} à jour (§\ref{sec:ngrok}). Un
        domaine statique supprime le problème.
  \item[\textmd{Pas de \texttt{Initialization Sequence Completed}, et le réseau filtre}]
        Le profil est en UDP 1194. Beaucoup de réseaux d'école ou d'entreprise
        ne laissent passer que TCP 443. Regénérer un profil TCP
        (§\ref{sec:openvpn-install}) et le redéposer (§\ref{sec:depot}).
  \item[\textmd{Tunnel « connecté », mais plus rien ne répond --- même Internet}]
        Deux VPN actifs en même temps. Couper l'autre
        (§\ref{sec:windows}).
  \item[\textmd{« Machine introuvable chez l'hébergeur »}]
        Nom, zone ou projet erroné, ou le compte de service n'a pas le rôle
        \code{compute.instanceAdmin.v1} (§\ref{sec:sa}).
  \item[\textmd{DVWA rame ou plante}] Swap absent : \texttt{free -h} (§\ref{sec:swap}).
\end{description}

\section{Résultat attendu}

\begin{itemize}[label=$\checkmark$]
    \item Projet GCP créé, facturation liée, API Compute activée
    \item \textbf{IP externe statique réservée} et attachée à la VM
    \item VM \texttt{e2-micro} dans une zone de la couche gratuite
    \item Pare-feu : UDP 1194 et TCP 22 seulement
    \item OpenVPN configuré avec l'IP \textbf{publique statique}
    \item Swap de 2~Go actif
    \item DVWA (\texttt{ghcr.io/digininja/dvwa}) en marche, port 80 restreint à \texttt{tun0}
    \item Tunnel monté, \texttt{ping 10.8.0.1} qui répond
    \item DVWA joignable à \texttt{http://10.8.0.1}, et \textbf{invisible} depuis Internet
    \item Compte de service GCP avec \code{compute.instanceAdmin.v1}
    \item \texttt{.ovpn} déposé dans la plateforme, téléchargeable par les apprenants
    \item Bouton « Démarrer / Arrêter » qui pilote réellement la VM
    \item Adresse \texttt{10.8.0.1} affichée par la plateforme une fois la machine en marche
    \item Plateforme exposée par ngrok, \texttt{APP\_PUBLIC\_URL} et
          \texttt{APP\_COOKIE\_SECURE} ajustés
\end{itemize}

\end{document}
